% AY2017 材料力学 〔II〕（転記）
% 出典: 04_材料力学/original/04_材料力学/2017/2017za.pdf
% 要約: 直角に折れ曲がった円形断面の部材（直径d）。固定端A → 区間AB(長さL) → 直角に
%       折れ → 区間BC(長さL) → 自由端C。先端C点に鉛直下向き外力P。ヤング率E, 横弾性係数G,
%       L >> d。
%       (1) B点ねじれ角ΨB, (2) B,C点のたわみ vB,vC,
%       (3) 固定端A断面の上表面の点D・中立面の点Eの垂直/せん断応力, モール円, 最大せん断/主応力。
% rem: 原本は3次元俯瞰図。直交2区間のL字部材を簡略立体図で再現する。

\section*{〔II〕}

図2に示すような直角に折れ曲がった円形断面の部材（直径 $d$）がある.
はりの先端（点 C）に垂直方向に外力 $P$ が作用するとき, 次の問い (1)\,$\sim$\,(3) に答えよ.
ただし, はりのヤング率 $E$, 横弾性係数 $G$ とする. はりの長さ $L$ は直径 $d$ に対して
十分大きいものとする.

\begin{center}
\begin{tikzpicture}[>=Latex]
  % 部材（L字, 円形断面を太線+丸端で表現）
  \draw[line width=14pt,gray!30,line cap=round] (1.9,1.7) -- (4.3,0.4) -- (7.3,1.4);
  \draw[line width=14pt,draw=black,fill=none,line cap=round,opacity=0] (1.9,1.7) -- (4.3,0.4) -- (7.3,1.4);
  % 端面（A, C）
  \draw[fill=gray!30] (1.9,1.7) ellipse (0.12 and 0.32);
  \draw[fill=gray!30] (7.3,1.4) ellipse (0.12 and 0.32);
  % 固定壁（A）
  \fill[pattern=north east lines] (1.35,1.0) rectangle (2.15,2.4);
  \draw (1.35,1.0) rectangle (2.15,2.4);
  \node[left] at (1.3,1.7) {\footnotesize 固定端};
  % 節点
  \node at (2.05,1.55) {A}; \node[below] at (4.35,0.3) {B}; \node[right] at (7.4,1.4) {C};
  % 軸（B付近）
  \draw[->] (5.0,0.7) -- (5.0,2.0) node[left]{$y$};
  \draw[->] (5.0,0.7) -- (3.7,0.95) node[above]{$z$};
  \draw[->] (5.0,0.7) -- (5.9,0.2) node[below]{$x$};
  % 外力 P（C点 下向き）
  \draw[->,very thick] (7.3,2.6) -- (7.3,1.55) node[above]{$P$};
  % 寸法 L（AB, BC）
  \draw[<->] (2.0,0.95) -- (4.3,-0.15) node[midway,below,fill=white,inner sep=1pt]{$L$};
  \draw[<->] (4.5,-0.05) -- (7.4,0.85) node[midway,below,fill=white,inner sep=1pt]{$L$};
  % A断面インセット（左下）
  \begin{scope}[shift={(0.6,-1.6)}]
    \draw[dashed] (0,0) circle (0.5);
    \draw[->] (0,0) -- (0,0.95) node[right]{$y$};
    \draw[->] (0,0) -- (-1.0,0) node[above]{$z$};
    \fill (0,0.5) circle (1.6pt); \node[right] at (0.05,0.55) {D};
    \fill (-0.5,0) circle (1.6pt); \node[above left] at (-0.5,0.05) {E};
    \node at (0,-0.95) {A 断面};
  \end{scope}
\end{tikzpicture}

\vspace{1mm}
図2
\end{center}

\begin{enumerate}[label=(\arabic*)]
  \item B 点でのねじれ角 $\Psi_B$ を求めよ.

  \item B 点および C 点でのたわみ $v_B$ および $v_C$ を求めよ.

  \item 固定端 A のはりの断面を考える. 図2に示す A 断面の上表面の点 D および中立面の点 E に
        生じる垂直応力およびせん断応力を求め, 点 D および点 E のモールの応力円を示せ.
        点 D および点 E の最大せん断応力および最大主応力を求めよ.
\end{enumerate}
